\scanpage{36}
is however possible that the injectivity of (*) remains valid without a hypothesis on the characteristic -- to be examined more closely...

Returning to the conditions of corollary 1, b), in characteristic $p>0$, and working with the prime-to-$p$ fundamental groups, one finds that the $u$ in 1°) is a point \inserted{of $G(K)$} of order a power of $p$. If $G$ is a torus (it suffices for \unsure{this|that one} that the ``fibers'' in the elementary variety $Y$ be \unsure{of|of} genus 0) one still concludes from this that $u=0$, i.e. $f=g$. In the general case, the arithmetic proof proves \inserted{(unless I am mistaken)} this : if $x\in X$, \deleted{and} \inserted{whence} $f(x)\inserted{=P}$, $g(x)=Q$ in $Y$, not only is $Q-P$ an element of $p$-torsion of the generalized \unsure{Jacobian|Jacobian} $G_Y$ of $Y$, for every finite Galois \inserted{étale} covering $Y'$ of $X$, \inserted{of degree prime to $p$\unsure{,|,,}} pointed by a $P'$ over $P$, there exists $Q'$ over $Q$ such that $Q'-P'$ is an element of ? $p$-torsion of $G_{Y'}$.

\deleted{by \unsure{taking|passing} as \unsure{coverings|liftings} the \inserted{abelian coverings, even \unsure{induced} by the} isogenies $\unsure{n_G|{}_nG|nG}$ ($n$ prime to \unsure{$G$|$p$}), there exists a $Q_n\in G$ such that $\unsure{{}_n n\cdot Q_n|n\cdot Q_n}=Q$, and that $Q_n$}
\unsure{It|I} could be that one could deduce $P=Q$ from this

\margin{This shows that \unsure{theorem|theorem} 2 remains true, over a base field of characteristic $p>0$, working with the prime-to-$p$ fundamental groups -- for there in the reduction to a geometric statement, one sees $X$, so that one can assume that $Y$ is a \uline{product} of curves of genus zero. In this case moreover, one is reduced to the case where $Y$ is itself such a curve,...}
\margin{35}
